\documentclass[a4paper,11pt]{article}
\usepackage[utf8]{inputenc}
\usepackage[T1]{fontenc}
\usepackage[magyar]{babel}
\usepackage{amsmath,amssymb}
\usepackage{geometry}
\usepackage{tcolorbox}
\usepackage{xcolor}
\usepackage{tikz}

\geometry{top=2cm, bottom=2cm, left=2cm, right=2cm}

\title{Hatványozás azonosságai_Dolgozat}
\author{AI LaTeX Studio}
\date{\today}

%%% Megjelenítéshez automatikusan kiegészített csomagok %%%
\usepackage{amsfonts, amsthm}
\usepackage{pgfplots}
\pgfplotsset{compat=1.18}
\usepackage{graphicx}
\usepackage{enumitem}
\usepackage{booktabs}
\usepackage{float}
\usepackage{caption}
\usepackage{microtype}

\usepackage{hyperref}

\begin{document}

% ==================== A CSOPORT ====================

\noindent
\begin{minipage}[t]{0.45\textwidth}
  \vspace{0pt}
  
  {\Large\textbf{Hatványozás azonosságai}}\\[2mm]
  {\large\textbf{„A” CSOPORT}}\\[2mm]
  \textbf{Név:} \underline{\hspace{4.5cm}}\\[2mm]
  \textbf{Osztály:} \underline{\hspace{1.5cm}} \quad \textbf{Dátum:} \underline{\hspace{2cm}}
\end{minipage}
\hfill
\begin{minipage}[t]{0.53\textwidth}
  \vspace{0pt}
  \begin{flushright}
  {\small
  \begin{tabular}{|l|c|c|c|c|c||c||c|}
    \hline
    \textbf{Feladat} & 1. & 2. & 3. & 4. & 5. & \textbf{Össz.} & \textbf{Érdemjegy} \\ \hline
    \textbf{Max. pont} & 8 & 10 & 8 & 10 & 8 & 44 & \\ \cline{1-7}
    \textbf{Elért pont} &  &  &  &  &  &  & \\ \hline
  \end{tabular}
  }
  \end{flushright}
\end{minipage}
\vspace{2mm}
\hrule
\vspace{4mm}
\begin{enumerate}
  \item \textbf{Alapműveletek és azonosságok alkalmazása} \hfill \textbf{[8 pont]}
  Add meg a következő kifejezések egyszerűsített alakját, és számítsd ki a helyettesítési értéküket, ha $x = 2$ és $y = -1$:
  \begin{enumerate}
    \item $\frac{2^3 \cdot 2^{-4}}{2^{-2}}$ \hfill \textit{[3 pont]}
    \item $(x^3 \cdot y^{-2})^2 \cdot \frac{x^{-1}}{y^3}$ \hfill \textit{[5 pont]}
  \end{enumerate}

  \item \textbf{Hatványok összehasonlítása és műveletek normálalakkal} \hfill \textbf{[10 pont]}
  \begin{enumerate}
    \item Rendezd növekvő sorrendbe az alábbi számokat anélkül, hogy kiszámítanád a pontos decimális alakjukat (írd le a gondolatmenetet): 
    $A = 2^{40}$, $B = 3^{30}$, $C = 5^{20}$. \hfill \textit{[4 pont]}
    \item Végezd el az alábbi műveletet, és az eredményt add meg normálalakban: 
    $\frac{4,8 \cdot 10^7 \cdot 3 \cdot 10^{-3}}{1,6 \cdot 10^2}$. \hfill \textit{[6 pont]}
  \end{enumerate}

  \item \textbf{Algebrai kifejezések és törtek egyszerűsítése} \hfill \textbf{[8 pont]}
  Egyszerűsítsd az alábbi törtet, és add meg a feltételeket is a változók lehetséges értékeire:
  $$\frac{a^3 b^{-2} \cdot (a^{-1} b^3)^2}{(a^2 b)^{-1}}$$
  feltéve, hogy $a \neq 0$ és $b \neq 0$.

  \item \textbf{Egyenletek megoldása hatványos alakban} \hfill \textbf{[10 pont]}
  Oldd meg a következő egyenleteket a valós számok halmazán:
  \begin{enumerate}
    \item $3^{x+1} + 3^x = 108$ \hfill \textit{[5 pont]}
    \item $2^{2x} - 5 \cdot 2^x + 4 = 0$ \hfill \textit{[5 pont]}
  \end{enumerate}

  \item \textbf{Gondolkodtató / Alkalmazási feladat} \hfill \textbf{[8 pont]}
  Egy biológiai kultúrában a baktériumok száma óránként megduplázódik. Kezdetben (0. órában) $2^{10}$ darab baktérium van a táptalajon. 
  \begin{enumerate}
    \item Mennyi baktérium lesz a kultúrában 5 óra elteltével? A választ írd fel egyetlen 2-es alapú hatványként! \hfill \textit{[3 pont]}
    \item Hány óra elteltével éri el vagy haladja meg a baktériumok száma a $2^{25}$ darabot? \hfill \textit{[5 pont]}
  \end{enumerate}
\end{enumerate}

\vfill
\noindent
\begin{center}
\small
\textbf{Ponthatárok:} \quad 1 (elégtelen): 0--17 p \quad 2 (elégséges): 18--23 p \quad 3 (közepes): 24--30 p \quad 4 (jó): 31--36 p \quad 5 (jeles): 37--44 p
\end{center}

\newpage

% ==================== B CSOPORT ====================

\noindent
\begin{minipage}[t]{0.45\textwidth}
  \vspace{0pt}
  
  {\Large\textbf{Hatványozás azonosságai}}\\[2mm]
  {\large\textbf{„B” CSOPORT}}\\[2mm]
  \textbf{Név:} \underline{\hspace{4.5cm}}\\[2mm]
  \textbf{Osztály:} \underline{\hspace{1.5cm}} \quad \textbf{Dátum:} \underline{\hspace{2cm}}
\end{minipage}
\hfill
\begin{minipage}[t]{0.53\textwidth}
  \vspace{0pt}
  \begin{flushright}
  {\small
  \begin{tabular}{|l|c|c|c|c|c||c||c|}
    \hline
    \textbf{Feladat} & 1. & 2. & 3. & 4. & 5. & \textbf{Össz.} & \textbf{Érdemjegy} \\ \hline
    \textbf{Max. pont} & 8 & 10 & 8 & 10 & 8 & 44 & \\ \cline{1-7}
    \textbf{Elért pont} &  &  &  &  &  &  & \\ \hline
  \end{tabular}
  }
  \end{flushright}
\end{minipage}
\vspace{2mm}
\hrule
\vspace{4mm}

\begin{enumerate}
  \item \textbf{Alapműveletek és azonosságok alkalmazása} \hfill \textbf{[8 pont]}
  Add meg a következő kifejezések egyszerűsített alakját, és számítsd ki a helyettesítési értéküket, ha $x = 3$ és $y = -2$:
  \begin{enumerate}
    \item $\frac{3^4 \cdot 3^{-2}}{3^{-1}}$ \hfill \textit{[3 pont]}
    \item $(x^{-2} \cdot y^3)^3 \cdot \frac{x^2}{y^{-2}}$ \hfill \textit{[5 pont]}
  \end{enumerate}

  \item \textbf{Hatványok összehasonlítása és műveletek normálalakkal} \hfill \textbf{[10 pont]}
  \begin{enumerate}
    \item Rendezd növekvő sorrendbe az alábbi számokat anélkül, hogy kiszámítanád a pontos decimális alakjukat (írd le a gondolatmenetet): 
    $A = 2^{30}$, $B = 3^{20}$, $C = 5^{10}$. \hfill \textit{[4 pont]}
    \item Végezd el az alábbi műveletet, és az eredményt add meg normálalakban: 
    $\frac{6,4 \cdot 10^8 \cdot 2 \cdot 10^{-4}}{1,6 \cdot 10^3}$. \hfill \textit{[6 pont]}
  \end{enumerate}

  \item \textbf{Algebrai kifejezések és törtek egyszerűsítése} \hfill \textbf{[8 pont]}
  Egyszerűsítsd az alábbi törtet, és add meg a feltételeket is a változók lehetséges értékeire:
  $$\frac{x^{-2} y^3 \cdot (x^2 y^{-1})^2}{(x^3 y)^{-1}}$$
  feltéve, hogy $x \neq 0$ és $y \neq 0$.

  \item \textbf{Egyenletek megoldása hatványos alakban} \hfill \textbf{[10 pont]}
  Oldd meg a következő egyenleteket a valós számok halmazán:
  \begin{enumerate}
    \item $2^{x+2} + 2^x = 320$ \hfill \textit{[5 pont]}
    \item $3^{2x} - 10 \cdot 3^x + 9 = 0$ \hfill \textit{[5 pont]}
  \end{enumerate}

  \item \textbf{Gondolkodtató / Alkalmazási feladat} \hfill \textbf{[8 pont]}
  Egy baktériumkultúrában a baktériumok száma óránként megháromszorozódik. Kezdetben (0. órában) $3^8$ darab baktérium van a táptalajon. 
  \begin{enumerate}
    \item Mennyi baktérium lesz a kultúrában 4 óra elteltével? A választ írd fel egyetlen 3-as alapú hatványként! \hfill \textit{[3 pont]}
    \item Hány óra elteltével éri el vagy haladja meg a baktériumok száma a $3^{20}$ darabot? \hfill \textit{[5 pont]}
  \end{enumerate}
\end{enumerate}

\vfill
\noindent
\begin{center}
\small
\textbf{Ponthatárok:} \quad 1 (elégtelen): 0--17 p \quad 2 (elégséges): 18--23 p \quad 3 (közepes): 24--30 p \quad 4 (jó): 31--36 p \quad 5 (jeles): 37--44 p
\end{center}

\end{document}